%% ============================================================
%% glossary.tex — Key Terms and Definitions
%% Process Intelligence Research Foundry
%% ============================================================
%%
%% This file defines canonical terms used throughout the dissertation.
%% Terms are defined using \newcommand for inline use and also
%% listed in the Glossary chapter for reader reference.
%% ============================================================

%% --- Inline term shortcuts ---
\newcommand{\ggen}{\texttt{ggen}}
\newcommand{\OCEL}{\textsc{ocel}}
\newcommand{\CodeManufactory}{\textsc{CodeManufactory}}
\newcommand{\ChatmanEq}{$\mathcal{A} = \mu(\mathcal{O}^*)$}
\newcommand{\ALIVE}{\textbf{ALIVE}}
\newcommand{\PARTIAL}{\textbf{PARTIAL}}
\newcommand{\BLOCKED}{\textbf{BLOCKED}}
\newcommand{\receipt}{\textit{receipt}}
\newcommand{\Receipt}{\textit{Receipt}}
\newcommand{\manufactured}{\textit{manufactured}}
\newcommand{\Manufactured}{\textit{Manufactured}}

%% ============================================================
%% Glossary Entry Definitions
%% (Rendered in the Glossary chapter or appendix)
%% ============================================================

%% To render the glossary in the document, include:
%%   \chapter*{Glossary}
%%   \addcontentsline{toc}{chapter}{Glossary}
%%   \input{glossary_entries}
%% or inline below if not using a dedicated glossary package.

\newcommand{\printglossary}{%
\chapter*{Glossary}
\addcontentsline{toc}{chapter}{Glossary}
\begin{description}

  \item[\textbf{ALIVE}]
    Status verdict indicating that all proof gates have passed for a given
    manufacturing artifact. Requires: receipt chain intact, all tests passing,
    checkpoint registered, and PDF hash recorded. An ALIVE verdict is permanent
    and immutable once issued.

  \item[\textbf{Atlas}]
    The canonical index of capabilities, algorithms, and execution surfaces
    across all research projects in the Process Intelligence Research Foundry.
    Not a summary --- a structured registry with evidence citations.

  \item[\textbf{BLOCKED}]
    Status verdict indicating that one or more core proof gates have failed.
    A BLOCKED verdict requires a remediation plan before any downstream
    manufacturing may proceed.

  \item[\textbf{Chatman Equation}]
    $\mathcal{A} = \mu(\mathcal{O}^*)$\quad where $\mathcal{A}$ is an autonomous
    agent, $\mu$ is the manufacturing operator, and $\mathcal{O}^*$ is the
    full object-centric process state. The equation asserts that a capable
    autonomous agent is produced by applying a deterministic, receipted
    manufacturing pipeline to a complete object-centric event log. Autonomy
    without process evidence is not autonomy --- it is hallucination.

  \item[\textbf{CodeManufactory}]
    The product. A deterministic, receipt-backed pipeline for manufacturing
    code and process artifacts from declared specifications. Not a framework,
    not a simulation, not a runner. RevOps is merely proof that CodeManufactory
    works.

  \item[\textbf{Conformance Checking}]
    The process-mining discipline of comparing a discovered process model
    derived from an event log against a declared normative model. Deviations
    are defects, not discrepancies.

  \item[\textbf{ggen}]
    Graph-query-template manufacturing engine. A Rust-based tool that manufactures
    structured artifacts (SPARQL queries, TOML configurations, YAML receipts,
    TeX chapters) from declared templates and ontology graphs. \ggen{} is
    itself a manufactured artifact and must carry a receipt chain.

  \item[\textbf{Manufactured}]
    Produced by a deterministic, receipted manufacturing pipeline. An artifact
    is manufactured if and only if: (1) its production was declared before
    execution, (2) every operator in the pipeline emitted a receipt, and
    (3) the receipt chain is verifiable. Artifacts produced outside this
    definition are not manufactured --- they are generated or hallucinated.

  \item[\textbf{OCEL}]
    Object-Centric Event Log. A process log format defined by Van der Aalst
    et al. (OCEL 2.0) that associates each event with multiple typed objects
    rather than a single case identifier. OCEL is the canonical evidence
    substrate for full-lifecycle process intelligence in this dissertation.

  \item[\textbf{PARTIAL}]
    Status verdict indicating that some but not all proof gates have passed.
    A PARTIAL verdict documents which gates passed and which remain open.
    Downstream manufacturing that depends on a PARTIAL artifact must declare
    the dependency explicitly and carry the PARTIAL risk in its own receipt.

  \item[\textbf{Post-Cyberpunk PCP}]
    Post-Cyberpunk Process-Complete Program. A software system that delivers
    working infrastructure with receipt-backed proof, not hallucination-as-output.
    The term distinguishes systems whose outputs can be replayed and verified
    from systems that produce plausible-seeming outputs with no evidentiary
    grounding.

  \item[\textbf{Process Evidence}]
    Event logs that prove a lawful process occurred. Not telemetry, not metrics,
    not dashboards. Process evidence is an OCEL-structured log from which a
    conformant process model can be discovered and verified against the declared
    normative model. If the log cannot support discovery, it is not evidence.

  \item[\textbf{Receipt}]
    A cryptographic or structured proof that a specific consequence occurred
    at a specific time in a specific manufacturing context. A receipt is
    not a log entry, a status flag, or a test result. A receipt is immutable,
    carries its own hash, and is chained to the receipt that preceded it.
    The receipt chain is the ground truth of the manufacturing pipeline.

  \item[\textbf{Replay}]
    Re-enacting a process from its event log. In process mining, replay
    is the mechanism by which a token or marking is moved through a declared
    Petri net model driven by the sequence of events in an observed log.
    Replay failures are conformance defects. In the CodeManufactory context,
    replay is the primary verification primitive: if the manufacturing pipeline
    cannot be replayed from its receipts, the pipeline is not verified.

  \item[\textbf{Witness Lattice}]
    A partially ordered set of type-law compatibility witnesses used in
    wasm4pm-compat analysis. Each witness asserts that a source type (e.g.,
    a pm4py EventLog) can be lawfully transformed into a target type (e.g.,
    a wasm4pm-compatible EventLog) under a specified set of structural
    constraints. The lattice orders witnesses by specificity of the
    transformation claim.

\end{description}
}
